\begin{saresult}{A.122}{Proved under stated hypotheses}{The complete component presentation in every PL dimension at least six}
Let $M$ be a compact connected PL manifold of dimension $n\geq6$,
with any boundary fixed on the sides of a concordance, and let $K$
be a finite simplicial model of its homotopy type. Then
\[
                   \pi_0P_{\rm PL}(M)\cong E(K),
\]
where $E(K)$ is the complete countable group presentation of
A.119. The same formula holds for the stipulated definable
concordance group of a compact definable manifold with this PL
model, over every real closed field and in every fixed o-minimal
expansion. Use the product of these presentations for finitely many
connected components, each of dimension at least six.

More precisely, positive suspension induces an isomorphism on
components of smooth concordance groups once $n\geq6$. This
identifies the unstable group of a smoothable PL model with the
stable presentation. The nonsmoothable PL case is obtained by
localizing to a smoothable neighborhood of a $3$-skeleton, not by
assuming a general unstable PL stability theorem. No assertion
about arbitrary manifolds of dimensions four or five is included.
\end{saresult}
\begin{proof}
We first lower the dimension bound for the smoothable case, keeping
track of the full relation subgroup in the corrected obstruction
sequence. For a connected space $X$ write
\[
 A_X=\operatorname{Wh}_1^+(\pi_1X;\mathbb Z/2\oplus\pi_2X),
 \qquad
 I_X=\operatorname{im}\bigl(\chi_X:
       \operatorname{Wh}_3(\pi_1X)\longrightarrow A_X\bigr).
\]
Replacing $\operatorname{Wh}_3$ by its surjecting algebraic
$K_3(\mathbb Z[\pi_1X])$ does not change $I_X$.
We use the corrected geometric exact-sequence theorem, in its full
form, as stated in
\cite[Theorem~3.1, p.~6]{Igusa2021Dimension3}: for a compact smooth
manifold of dimension at least six, it gives
\[
 0\longrightarrow A_M/I_M\xrightarrow{\ \overline\theta_M\ }
 \pi_0P_{\rm Diff}(M)\xrightarrow{\ \Sigma_M\ }
 \operatorname{Wh}_2(\pi_1M)\longrightarrow0.             \tag{A.122.1}
\]
The $\chi$ here is the full corrected obstruction map, including
its framing component. It is the same map as in the stable
corrected sequence in \cite[Introduction, (0.2), p.~2]
{Igusa2021SecondObstructionI}, and in the functorial A-theory
construction of \cite[Section~2, especially Corollary~2.7, p.~250]
{Jahren2010}. This identification of the full corrected sequence is
a cited input. In particular, we are not deriving it from the
formula for the $\pi_2M$-component of $\chi$ printed after Igusa's
Theorem~3.1. That partial formula would not determine $I_M$.
The original 1984 chapter \cite{Igusa1984Postnikov} has not been
retrieved; the precise opened restatements and functorial
construction just listed are the sources used here.

The deformation retraction $M\times I\longrightarrow M$
identifies $\pi_1$, its action on $\pi_2$, and their first Postnikov
invariant. Functoriality of the full sequence in
\cite[Corollary~2.7]{Jahren2010} identifies $A_{M\times I}$ with
$A_M$ and $I_{M\times I}$ with $I_M$. We next verify that these
are the vertical identifications induced by positive suspension on
the two geometric obstruction maps in (A.122.1).

Use concordances that are the identity near the bottom and sides,
and are products near the top. Inclusion of this normalized
subgroup induces the same component group; collar straightening
also works for compact one-parameter families. The positive
suspension formula of
\cite[Section~4.1, p.~8]{Igusa2021Dimension3} applies a
concordance $g$ on each fixed polar ray and is the identity on the
remaining region. Thus
\[
                     \sigma^+(gh)=\sigma^+(g)\sigma^+(h)
\]
on that subgroup. It consequently gives a homomorphism of
component groups, respecting the bottom and side collars.

For the second obstruction use marked lens models in the fixed
indices $3,4$ in every dimension $n\geq6$. This convention avoids
an index change when $n$ changes parity. The two-index deformation
is available in this range: the original
\cite[Part~I, Chapter~V, Theorem~3.1(a), p.~184]
{HatcherWagoner1973} permits $2\leq i\leq n-2$, and its use for
the one-parameter representatives is explicit in Part~II, Section~2,
p.~244. The marked-lens framed-bordism calculation of
\cite[Theorem~1.8 and Lemma~1.10, pp.~8--9]
{Igusa2021SecondObstructionI} applies with $k=3$: its proof needs
the inequalities $k\geq3$ and $n-k\geq3$, which hold here.
It identifies the marked-lens group with
\[
 (\mathbb Z/2\oplus\pi_2M)[\pi_1M],/
       (\mathbb Z/2\oplus\pi_2M)[1].
\]
Changing the markings gives the coinvariants $A_M$; forgetting the
lens model gives $\theta_M$. Use the stable sign convention
$\lambda_0^s=(-1)^k\lambda_0$ from
\cite[Section~4.2, pp.~8--9]{Igusa2021Dimension3}.
That source explicitly proves compatibility with positive
suspension, which retains $k$, for the complete framed invariant.
This includes its $\mathbb Z/2$ coordinate. Hence the square
\[
 \sigma^+_*\theta_M=\theta_{M\times I}
       : A_M\longrightarrow\pi_0P_{\rm Diff}(M\times I)
                                                               \tag{A.122.2}
\]
commutes under the preceding identification of $A$.
The first obstruction is unchanged under the same suspension:
\cite[Part~I, Chapter~IV, Section~1, Remark, p.~130]
{HatcherWagoner1973}, using the trajectory comparison in
Chapter~I, Section~5, (5.2), p.~60. Therefore
\[
       \Sigma_{M\times I}\sigma^+_* =\Sigma_M.          \tag{A.122.3}
\]

Equations (A.122.2)--(A.122.3) give a commutative diagram of
short exact sequences (A.122.1), whose left and right vertical
arrows are the identity on $A_M/I_M$ and
$\operatorname{Wh}_2(\pi_1M)$, respectively. The middle arrow is
$\sigma^+_*$. It is an isomorphism by exactness. Explicitly, if
$\sigma^+_*(g)=1$, then $\Sigma_M(g)=0$, so
$g=\overline\theta_M(a)$; injectivity of the lower
$\overline\theta$ forces $a=0$ and hence $g=1$.
Given $g'$ in the target, first choose $h$ with
$\Sigma_M(h)=\Sigma_{M\times I}(g')$. The difference
$g'\sigma^+_*(h)^{-1}$ lies in the image of
$\overline\theta_{M\times I}$ and is therefore the suspension of
an element of $\operatorname{im}\overline\theta_M$. This proves
surjectivity. This argument does not split either extension.

The compact-manifold and boundary hypotheses persist throughout.
In \cite[Introduction, p.~1]{Igusa2021SecondObstructionI},
$P_{\rm Diff}(M)$ is explicitly
$\operatorname{Diff}(M\times I\text{ rel }M\times0\cup
\partial M\times I)$. Section~1, p.~3, fixes the projection
function near the entire boundary. The lens modifications and
two-index deformations are supported in the interior. The
Hatcher--Wagoner setup likewise uses compact pairs $(M,\partial M)$,
and the collar normalization preceding suspension is on printed
p.~55. Thus the argument applies to every $M\times I^j$ after
rounding its corners. Iterating the component isomorphism
identifies $\pi_0P_{\rm Diff}(M)$ with its stable group once
$n\geq6$. The complete stable group presentation in A.119 is
$E(K)$, with its group law already checked through classifying
h-cobordisms. The smooth-to-PL component isomorphism of A.103
therefore proves the desired statement whenever $M$ is smoothable
and $n\geq6$.

Now let $M$ be any compact PL manifold of dimension $n\geq7$.
Repeat the finite regular-neighborhood construction of A.121.
All of its localization hypotheses require only $3\leq n-4$ and
$2+3<n$, which hold here. The resulting $N$ is a smoothable
$n$-manifold with the homotopy type of the chosen $3$-skeleton.
The just-proved smoothable result now applies to $N$ in this lower
dimension. Extension by the identity consequently gives
\[
                 \pi_0P_{\rm PL}(M)\cong E(K^{(3)}).
                                                               \tag{A.122.4}
\]
Here first take $K$ to be a triangulation of $M$.
There is a natural group isomorphism
$E(K^{(3)})\longrightarrow E(K)$. To see it, use the functorial
sequence of \cite[Corollary~2.7]{Jahren2010} for the inclusion
$K^{(3)}\longrightarrow K$. This inclusion is an isomorphism on
$\pi_1$ and $\pi_2$, so it induces isomorphisms on
$\operatorname{Wh}_3$, $A$, and $\operatorname{Wh}_2$.
The square with the full $\chi$ commutes, so it also induces an
isomorphism on its cokernel. Applying exactness to
\[
 0\longrightarrow A_X/I_X\longrightarrow
 \pi_2\operatorname{Wh}^{\rm Diff}(X)\longrightarrow
 \operatorname{Wh}_2(\pi_1X)\longrightarrow0
\]
therefore gives an isomorphism on the middle groups. A.119 identifies
these groups with $E(K^{(3)})$ and $E(K)$, proving the claim.
This also allows any finite simplicial homotopy model $K$ in the
statement, rather than only the selected triangulation.

For completeness, the remaining dimension six case needs no
localization: every PL $6$-manifold admits a compatible smoothing.
Here is the exact classical input and its application. A finite PL
handle decomposition has handle indices at most six. The
PL--DIFF extension argument in
\cite[Essay~IV, Appendix~B, p.~211]{KirbySiebenmann1977}
identifies the obstruction at a $k$-handle with the class of a
smooth twisted $(k-1)$-sphere, namely an element of
$\Gamma_{k-1}$. These groups vanish for $0\leq k-1\leq5$.
For the low-dimensional compatible-smoothing assertion one can
read \cite[Essay~I, p.~39]{KirbySiebenmann1977}; the remaining
values $\Gamma_4=0$ and $\Theta^{\rm Diff}_5=0$ are explicitly
stated in Essay~IV, Section~10.12, p.~200. A smooth twisted
$5$-sphere represents a smooth homotopy $5$-sphere, so the latter
equality gives $\Gamma_5=0$. Starting with the $0$-handles and
extending across each successive handle now produces a compatible
smooth structure on $M$, including a smooth structure on its
boundary. No preassigned boundary smoothing is required here.
The smoothable case already proved applies and gives $E(K)$.

Finally A.113 and A.115 transfer the PL component calculation to
the stipulated definable concordance group over any real closed
field. Their relative versions preserve the side-boundary
convention and composition. Finite products handle disconnected
manifolds as in A.121. The lower-dimensional exclusions in the
statement remain: neither the corrected exact-sequence input nor
the fixed-index argument used here asserts this conclusion for
arbitrary $4$- or $5$-manifolds.
\end{proof}
